% Supplementary Methods: the detail behind the main text's Methods (participants, materials and
% procedure, measures, registrations, models, model checks, the registered refits, the text
% measures, the exploratory analyses, software, random seeds and run records). It points to each
% study's section, to the Registrations section and to the Supplementary Notes rather than repeating
% them. Numbers come only from macros in generated/numbers.tex; design facts and analysis settings
% are typed. This section has no numbered table or figure, so that the Supplementary Table and
% Figure numbers of the sections after it do not change.
\section{Supplementary Methods}\label{sec:supp-methods}

The main text's Methods give the design and analysis of the five studies in outline. This section gives the detail behind them, in the same order. It does not repeat what other sections hold. Each study's section quotes that study's materials and gives its exclusions step by step and its sample. Section~\ref{sec:suppA-registrations} gives the registrations, what we had seen of each study's data before filing them, how the images were made, and the corrections to earlier reports. Section~\ref{sec:suppA-text} describes the word lists, and the Supplementary Notes report the results. A registration here is a dated record of a study's hypotheses and analysis plan, not a copyright registration.

\subsection*{Participants}

\paragraph*{Sample sizes.}
No power analysis, a calculation of how many participants a study needs to detect an effect of a given size, was run for any study. Each survey set a quota, a planned number of participants: 60 for each label group in Study~1, 125 in Study~2, 50 in Study~3 and 62 in Study~4, and 40 for each of Study~5's six combinations of label and editing time. No quota sent anyone to the other label. Study~2's planning document, which was never submitted, had set a target of 210 participants and a stop at 250. Study~1's registration adds simulations of what its design could detect, run after its data had been collected (Supplementary Note~1); they did not set its sample size. In Studies 2 to 5, a claim that an effect is small rests on the equivalence tests described below.

\paragraph*{Exclusions.}
Each registration states its exclusion rules and their order. The rules were written after the data had been collected, and earlier analyses of Studies 3 and 4 had kept some records that these rules remove. Every study first removed survey test runs and unfinished sessions. The closing screen asked `Did you complete this experiment seriously throughout (without responding randomly)?', and every study then removed anyone who gave any answer but `Yes, I did', or no answer. Studies 2 to 5 next removed wrong answers to an arithmetic question, `What does 8+7 equal?'. Last, Studies 2 and 5 removed records with no recorded condition (there were none in Study~2), and Study~4 removed records that were not valid completions. A valid completion showed exactly six originals, each once, two at each stated time. In Study~1 the sequence also includes the question about the task, which removed no one, because the survey let everyone continue after one retry. Each study's section gives the number removed at every step.

\paragraph*{Participations and people.}
The five analyzed samples add up to \NTotal{} participations, where a participation is one person's session in one study. We counted the people behind them by matching the subject-pool participant code, the code that the survey returned at its end to the pool's page for granting credit and that is the same for a person in every study. The \NTotal{} participations came from \NPeople{} people. Of these, \NPeopleRepeat{} took part twice, each in one earlier study and then in Study~5. Studies 1 to 4 ran under one subject-pool listing and Study~5 under a listing of its own, which would let earlier participants sign up again. A person in Study~5 who had taken part before had already seen some of the same originals and read the debriefing, which told them that anything produced with AI is insufficiently creative (below). No registered analysis removes them, no analysis checks whether their first session changed their ratings, and the exploratory pooled model of Studies 2, 4 and 5 counts twice those who had taken part in Study~2 or~4. This count, like the age and gender summaries below, was specified on 19 October 2026, after the registered results were known; it is descriptive and tests nothing.

\paragraph*{Age, gender and other background questions.}
The closing screen asked for race, gender, age, native language and formal artistic training. Age was typed as free text and read as the first number in the answer. Gender was chosen from four answers: male, female, other and prefer not to say. No question asked about sex. The table below gives, for each analyzed sample, the number of participants under each label and a summary of age and gender. `Other or none' joins three answers: other, prefer not to say and no answer.

\begin{center}
\small
\begin{tabular}{@{}lrrrrrrrr@{}}
\toprule
Study & $N$ & Human & AI & Mean age (SD) & Age range & Female & Male & Other or none\\
\midrule
Study~1 & \SOneN & \SOneNHuman & \SOneNAI & \SOneAgeMean{} (\SOneAgeSD) & \SOneAgeMin--\SOneAgeMax & \SOneNWomen & \SOneNMen & \SOneNGenderOther\\
Study~2 & \STwoN & \STwoNHuman & \STwoNAI & \STwoAgeMean{} (\STwoAgeSD) & \STwoAgeMin--\STwoAgeMax & \STwoNWomen & \STwoNMen & \STwoNGenderOther\\
Study~3 & \SThreeN & \SThreeNHuman & \SThreeNAI & \SThreeAgeMean{} (\SThreeAgeSD) & \SThreeAgeMin--\SThreeAgeMax & \SThreeNWomen & \SThreeNMen & \SThreeNGenderOther\\
Study~4 & \SFourN & \SFourNHuman & \SFourNAI & \SFourAgeMean{} (\SFourAgeSD) & \SFourAgeMin--\SFourAgeMax & \SFourNWomen & \SFourNMen & \SFourNGenderOther\\
Study~5 & \SFiveN & \SFiveNHuman & \SFiveNAI & \SFiveAgeMean{} (\SFiveAgeSD) & \SFiveAgeMin--\SFiveAgeMax & \SFiveNWomen & \SFiveNMen & \SFiveNGenderOther\\
\bottomrule
\end{tabular}
\end{center}

\noindent Human and AI give the number of participants under each label. SD is the standard deviation, a measure of how far ages typically lie from the mean. Supplementary Table~\ref{tab:supp_demographics} gives Study~1's answers on race and artistic training as well; the answers of Studies 2 to 5 to those questions are not summarized here. No analysis compared genders or other background groups. Age, gender and training entered only some of Study~1's robustness refits, as covariates, variables added to a model to adjust for them (below).

\subsection*{Materials and procedure}

\paragraph*{The label.}
Every version of the label opened `There is a digital artist named Alex. Alex is a skilled artist who creates digital artworks.' The AI label went on, in every study: `Alex develops and inputs prompts into an artificial intelligence (AI) system. The AI system generates digital artworks based on those prompts.' The Human label went on, in Studies 1, 4 and 5: `Alex uses his graphic design skills and digital sketching tools.' In Study~2 it went on `Alex only uses graphic design software and digital sketching tools', and in Study~3 `Alex uses only his graphic design skills and digital sketching tools such as iPad apps and Adobe Photoshop.' Before every original, a short description restated how Alex had made that image. These descriptions differed slightly across studies: in Study~2, the AI version said that Alex prompted the AI system without saying that he revised his prompts, and in Studies 3 to 5 both versions stated a time. In Study~5 the description of each second work stated the editing time in the same words at every level, so at None it said that Alex had spent that time on an unchanged copy. Each study's section quotes every description that participants read.

\paragraph*{The images.}
Section~\ref{sec:suppA-registrations} (How the images were made) describes how the ten originals were generated, the provenance records that eight of their files carry, the prompts and the hand edits. Supplementary Note~4 identifies the \XSImageCheckSlightReframedN{} Slight versions that were also cropped or resized. For that check, the 30 images were extracted, without being re-encoded, from a printout of the Study~5 survey: a PDF file printed from the survey software, which holds each image as the survey served it.

\paragraph*{Screens and order.}
Each survey opened with a welcome screen and then the consent form. The instructions said that participants would be asked to explain their answers. A question about the task followed (`What is your task in this experiment?'); a wrong or missing answer brought back the instructions and the question once, after which every participant went on. Each original stayed on screen alone for at least 8 seconds. In Studies 2, 4 and 5, each second work was introduced and then shown beside its original, headed `FIRST WORK' and `SECOND WORK', for at least 10 seconds. Each rating screen gave all of an image's ratings in a fixed order (creativity, then effort; or transformativeness, essence change, then creativity) and below them a required box for explaining each rating. Copying and pasting were blocked in none of the 20 explanation boxes of each of Studies 1 and 3, and in 21 of the 90 boxes of each of Studies 2, 4 and 5.

The originals came in a random order, except that in Studies 4 and 5 the three with a figure came together, as did the three without, the two groups in random order. The three second works of each original also came in a random order. The surveys did not record the order of the originals in Studies 1 and 2, or the order of the second works in any study. In Study~5, the survey script meant to record each original's source time looked for it under a field name that the survey no longer used, so that record is blank for every participant. The arithmetic question came at a random point among the originals in Study~2 and after the last rating in Studies 3 to 5; Study~1 had none. The closing questions followed in this order: the seriousness question, a question about problems with the procedure, the background questions and a box for any other comments.

\paragraph*{The debriefing and the law.}
The last screen, the same in all five studies, read: `We are studying how humans perceive the copyrightability of artificial intelligence art. Current legal standards state that anything produced with artificial intelligence is insufficiently creative, even if you use AI as a base against which you create new work. There is ongoing litigation against the US Copyright Office with an analogous set of facts that you have seen here.' A final sentence invited questions. The second sentence overstates the law. The US Copyright Office refuses to register material whose expressive elements an AI system determined. It does register a person's own creative selection or arrangement of that material, or a modification of it creative enough to be original, with the AI-generated material left out of the claim.\cite{usco2023guidance,usco2025copyrightability} The Office treats such a modification by analogy to a derivative work, a work based on an earlier work, whose copyright covers only the material its author contributed (17 U.S.C.~\S\,103(b)). It is only an analogy, because in the Office's view AI output is not itself a work of authorship.\cite{usco2025copyrightability} The debriefing instead stated as law the categorical view that this paper calls the original sin: that AI involvement leaves even work built on it insufficiently creative. The paper tests that view as a hypothesis about lay judgment, not as a legal rule. The debriefing also did not tell participants that every image had been made with AI.

\subsection*{Measures}

Each study's section quotes its rating items and answer labels. Studies 4 and 5 defined essence as `the core sense and feeling that one has while observing a work of art that makes that work unique', and defined the creativity of a second work as for originals. Study~2 defined neither. In each of Studies 2, 4 and 5, two transformativeness items had a sixth answer point, which puts them on a different scale. In Study~2 these two items also asked `How transformative do you consider the second work compared to the first work?', with no definition. As registered, their ratings are left out of the primary analyses: \STwoNSixPointDropped{} in Study~2, \SFourNSixPointDropped{} in Study~4 and \SFiveNSixPointDropped{} in Study~5. Study~5's registered second analysis puts them back on a common scale (below).

\subsection*{Registrations}

Section~\ref{sec:suppA-registrations} lists the six registration documents and, study by study, what we had seen of each study's data before filing. Two points from the Methods are given here in full. First, for Studies 2 to 5 the code reported here was run before filing only with the condition labels shuffled at random across participants, so that every model and table could be checked while no number described the real conditions. In Study~5's shuffled runs, and in one early shuffled run of Study~3's, the ratings were not shuffled, so the text analyses that relate the words to the rating alone used the real ratings. Those outputs were set aside, and none is reported. Second, the CogSci 2025 paper,\cite{kamper2025originalsin} which first reported Studies 2 and 4, makes several statements that this paper corrects. It called its first study, this paper's Study~2, pre-registered, but no study was registered anywhere before its data were collected. It attributed the prompts to a validated set published by Orwig and colleagues,\cite{orwig2026using} but we wrote the prompts ourselves. It described the Slight versions as adjustments of hue, coloration and contrast, without the cropping and resizing. It said that a power analysis had set its sample size, but none was run. It gave Study~4's sample as \SFourNBeforeValidity, the count before the registered rule that removes \SFourNInvalidRemoved{} invalid completions. Section~\ref{sec:suppA-registrations} sets out each correction, with the others.

\subsection*{The registered models}

\paragraph*{What every model shares.}
Each registered analysis is a linear mixed model, a regression that also gives each participant and each original its own baseline, called a random intercept. One model was fitted for each rating question, to every individual rating rather than to participants' averages. The models used \SOneNRatings{} ratings in Study~1, \STwoNRatings{} in Study~2, \SThreeNRatings{} in Study~3, \SFourNRatings{} in Study~4 and \SFiveNRatings{} in Study~5. Missing ratings were left out, not filled in. In Study~2, \STwoNLackingOneTrial{} participants lack one trial. Study~3's records lack one rating for each of \SThreeNLackingOneRating{} participants, and \SThreeNAbsentCode{} other ratings carry an answer code that the final survey did not offer, so they are treated as missing. Models were fitted by restricted maximum likelihood, the usual way of estimating such models, with the R packages lme4 and lmerTest. Studies 2 to 5 coded each factor so that its effects are measured from the average over its levels. Study~1 coded the label and the kind of image against reference levels (Human, and no figure), so its penalty is the averaged difference described next, not a single model coefficient.

Each reported difference is a contrast between estimated marginal means: the mean ratings that the model predicts for two conditions, each averaged with equal weight over the other factors. $F$~tests asked whether a factor changed the mean rating, and $t$~tests whether a single difference departed from zero. Their degrees of freedom, which set how large a statistic must be to count as significant, came from Satterthwaite's approximation, computed for every rating, including in the larger studies, where the software would otherwise have switched to a large-sample approximation. Study~1's contrasts used the Kenward--Roger approximation, as its long registration specifies; it gives the same value, \SOneHOneCreatDf, for the penalty. A $p$-value is the chance of a result at least this far from no effect if there were none. Tests were two-sided at .05, and a 95\% confidence interval (CI) is the range of values that the data do not rule out at that level.

\paragraph*{The model of each study.}
In the formulas below, a star joins factors whose effects and interactions (changes in one factor's effect across the levels of another) are all estimated, and a term such as $(1 \mid \textrm{participant})$ gives each participant a baseline of its own.
\begin{itemize}\raggedright
\item Study~1, per rating: rating $\sim$ label $*$ kind of image${}+ (1 \mid \textrm{participant}) + (1 \mid \textrm{original})$. H1 holds if the penalty on both creativity and effort is positive and survives Holm's correction. H2 is tested in one joint model of both ratings, rating $\sim$ label $*$ rating question $*$ kind of image${}+ (1 + \textrm{rating question} \mid \textrm{participant}) + (1 \mid \textrm{original}{:}\textrm{rating question})$, which lets each participant's gap between the two ratings differ. Its corrected test is the $F$~test of the label by rating question interaction, and an ordinal version of the same model must agree in sign and significance. H3 is the label by kind of image interaction. It is judged by a two-sided test together with an equivalence test against $\pm$\SOneHThreeBound, through a four-cell decision table in the long registration. A failed equivalence test is uninformative if the test's power, its chance of passing when the true difference is zero, is below 80\%.
\item Study~2, per measure: rating $\sim$ label $*$ modification level $*$ kind of image${}+ (1 \mid \textrm{participant}) + (1 \mid \textrm{original})$. Four tests per measure enter the corrected set of \STwoNConfirmatory: Slight minus None and Dramatic minus Slight, which together decide H1; the penalty averaged over levels, tested for equivalence (H2); and the penalty at None minus the penalty at Dramatic (H3).
\item Study~3, per rating: rating $\sim$ label $*$ stated time${}+ (1 \mid \textrm{participant}) + (1 \mid \textrm{original})$. The corrected set holds the $F$~tests for label, time and their interaction on each rating. The registered decision calls the gap closed by time only if the interaction survives Holm's correction, the gap is smaller at 100 hours than at under 1 hour, and the gap at 100 hours passes the equivalence test. It calls the gap persistent if that equivalence test fails and the gap's 95\% CI at 100 hours lies above zero, and otherwise inconclusive.
\item Study~4, per measure: rating $\sim$ label $*$ source time $*$ modification level${}+ (1 \mid \textrm{participant}) + (1 \mid \textrm{original})$. Five tests per measure enter the corrected set of \SFourNConfirmatory: the two steps of H1, the equivalence test of H2, and the $F$~tests for source time and for source time by modification level (H3). Registered secondary tests, judged at .05 without correction, give the penalty at each level, the label by source time interaction, and an equivalence test of 100 hours minus under 1 hour against $\pm$\SourceBound.
\item Study~5, per measure: rating $\sim$ label $*$ editing time $*$ modification level $+$ kind of image${}+ (1 \mid \textrm{participant}) + (1 \mid \textrm{original})$, with kind of image added on its own. The corrected set holds the $F$~tests for editing time and for label by editing time on each measure. Registered secondary equivalence tests, not corrected, ask whether 100 hours minus under 1 hour of editing time lies inside $\pm$\SFiveEditTimeBound{} and whether the penalty at each level lies inside $\pm$\SourceBound; the verdict on editing time also reads its corrected $F$~test.
\end{itemize}

\paragraph*{H4.}
H4 asked whether the AI label spreads ratings more between participants and less within each participant, across images. Under each label it estimated two variances, a variance being a measure of spread: the variance between participants, and the variance within each participant across images. The first came from a model that lets the variance between participants differ by label, and the second from a model that lets the remaining variance differ by label. It compared the labels by the ratio of each variance, AI over Human, with 95\% intervals from 1,000 resamples of participants within each label, so its result is a pair of ratios, not a difference between means. H4 holds on a rating only if the ratio between participants lies above 1 and the ratio within participants below 1, each with its whole interval on that side, and it must hold on both ratings. Its entry in Study~1's corrected set is a likelihood-ratio test, which asks whether a model that lets the variance between participants differ by label fits better than one that does not; the rating with the larger $p$-value enters. Supplementary Note~1 reports the result.

\paragraph*{Equivalence tests and bounds.}
An equivalence test is two one-sided tests, each at .05, of whether an effect lies above a lower bound and below an upper bound.\cite{schuirmann1987comparison,lakens2017equivalence} Its $p$-value is the larger of the two one-sided $p$-values, and judged alone it passes when the effect's 90\% CI lies wholly inside the bounds. In Studies 2 and 4, that $p$-value entered Holm's correction with the study's other confirmatory tests, so an interval can lie inside the bound while the corrected test fails. Study~1's bound of $\pm$\SOneHThreeBound{} was set after its analyses had been run, on the ground that a smaller difference is finer than the resolution at which participants use a five-point scale. The bound of $\pm$\SourceBound{} is half of Study~1's registered creativity penalty of \SOneHOneCreatDiffPrecise{} points, used for all three second-work measures for lack of another benchmark. Study~5's registration fixed the rule for its editing-time bound before any Study~5 model was fitted. Its code read the value, Study~3's registered creativity contrast between 100 hours and under 1 hour (\SThreeTimeContrastPrecise), from the output of Study~3's run with the real labels, and would have stopped had that output come from a shuffled run.

\paragraph*{Effect sizes.}
Study~1's registered standardized effect size is Cohen's $d$: the difference between the label groups' mean ratings divided by the spread of participants' mean ratings, with a 95\% interval from 10,000 resamples. A model-based version divides the penalty by the square root of the summed variances of the model, so that it can be compared with other studies. Section~\ref{sec:suppA-study1} reports both.

\subsection*{Ordinal checks and fitting status}

\paragraph*{Ordinal checks.}
As registered checks, the confirmatory models were refitted as cumulative-link mixed models, which treat the five answers as ordered categories rather than evenly spaced numbers and report estimates on the log-odds scale, the logarithm of the odds of a higher answer. They used the R package ordinal, with a logit link and the Laplace approximation, and their tests are Wald tests, which compare an estimate with its standard error. In Study~1, agreement of the ordinal model in sign and significance was part of the registered rule for H2. In Study~2, the nine ordinal $p$-values for the level steps and the dilution test replaced the linear ones in a corrected set of twelve. The three equivalence tests stayed as they were, because their bound is in rating points and has no ordinal counterpart. Study~4 did the same with its twelve tests of differences in its set of fifteen, and Studies 3 and 5 corrected the ordinal test of each model term over the same six tests. Elsewhere than Study~1's H2, the registrations call the ordinal model a check, so a disagreement is reported beside the verdict rather than replacing it.

\paragraph*{Fitting status.}
Every fitted model carries a status. A linear mixed model converged if its optimizer, the numerical search for the best-fitting values, reported success and raised no warning about convergence. An ordinal model converged if its optimizer reported success and raised no warning about convergence or about the shape of its fit. A model that did not converge stopped without settling on a solution. It yields no $p$-value, interval or verdict, and its check is reported as not completed. Study~1's registration added this rule in October 2026, after its September estimates had been read, and every study's code applies it. The fits that did not converge were the ordinal checks of Study~2's creativity model and of the pooled creativity model, and the structural topic model of the essence-change explanations. A singular fit converged but put a variance, or a correlation between baselines, at the edge of its possible range, such as a variance of zero, a sign that the data cannot support every quantity the model asks for. Such a fit keeps its numbers and is reported as singular. The singular fits include Study~1's most complex refits, Study~3's refits with random slopes and every pooled linear mixed model of the exploratory analyses.

\paragraph*{Random slopes.}
Two registrations asked that a verdict also hold in a model with random slopes, which let an effect differ from one participant, or one original, to the next. Study~1's asked this of H3 and Study~3's of its equivalence test. Each let the simpler model decide if the slope model was singular, and Study~3's also if that model stopped with an error or did not converge. Study~1's registration does not list this rule among the three elements fixed before its analyses were run. In both studies the slope model was singular, so the verdict rests on the simpler model.

\subsection*{Study~1's robustness checks and the other registered refits}

\paragraph*{Study~1.}
The registered robustness checks refit the penalty in several ways (Supplementary Note~2). A specification curve refit it under every combination of a set of defensible analysis choices. The choices were four. The model had baselines only, or the most complex structure of random slopes. It was a linear mixed model, an ordinal model, or generalized estimating equations, which estimate an average effect while allowing ratings from one participant to be related. It had no covariates, background covariates, or those plus the length of the explanations. And it was fitted to every rating, or to each participant's mean for each kind of image, with standard errors that allow for related ratings. The combinations reduce to \SOneSpecN{} distinct models for each rating. Other refits added background measures as covariates, left out one participant or one original at a time, dropped one effort item whose answer labels differed from the rest, or used Bayesian models.

Bayesian models combine the data with stated assumptions about plausible values, called priors. The four Bayesian refits, a linear and an ordinal version for each rating, ran in brms. Their priors were normal(0, 1) on the effects, a bell curve centered on zero with a standard deviation of 1 point, and Student-$t$(3, 0, 2.5) on the intercept and the standard deviations, a wider curve with heavier tails. Each model drew its estimates in four independent runs, called chains, of 4,000 iterations each. The first 2,000 of each run were discarded as warm-up, while the sampler tuned itself. The sampler took small steps (an adapt\_delta setting of 0.95) to avoid failed steps, called divergences. A fit counted as converged only if, for every quantity, the chains agreed (R-hat below 1.01), the draws were worth more than 400 independent ones, and none diverged.

\paragraph*{The belief refits and the rescaled analysis.}
A registered refit of Study~1 left out the Human-label participants who wrote at least once that a work is or looks machine-made, by the word list for a statement that a machine made the work. As exploratory description, 2,000 random groups of Human-label participants of the same size and mean effort rating were removed in the same way for comparison (Supplementary Note~2). Study~4's registration added a refit without the Human-label participants who wrote that a machine made the work, which repeats all fifteen corrected tests (Section~\ref{sec:suppB-study4}). Study~5's registration added a second analysis that rescales every item to run from 0 to 1, as (rating $-$ 1) divided by (number of answer points $-$ 1), so that the six-point items can enter. It repeats the six corrected tests and runs no equivalence test, because both bounds are in rating points.

\subsection*{The text measures}

Section~\ref{sec:suppA-text} describes the eight word lists, how they were written and applied, and the two registered checks. The lists flag visible detail, time, skill, concept (the idea behind the work), prompting or machines, a claim that the work looks familiar, the scale's own answer labels, and a statement that a machine made the work. The script that applies them received only identifiers and text, and stopped if handed the label, the rating or any background measure; the label and the rating were joined to the flags afterwards. The first registered check asks how well the words predict the rating and the label for participants left out of the fitting, using ridge regression, a regression that shrinks each word's weight toward zero. The second asks whether participants under the AI label reuse their own words across explanations more than those under the Human label, relative to other participants writing about the same picture. Study~1's registration also names five tests of what its effort and creativity explanations referred to. They form an exploratory set corrected by the Benjamini--Hochberg method, which limits the expected share of false positives among the results called significant, and their corrected results are descriptive readings, not verdicts.\cite{benjamini1995controlling} All other analyses of the explanations are exploratory, among them scores of word tone and concreteness from published norms, average ratings of single words collected by other researchers.\cite{warriner2013norms,brysbaert2014concreteness}

\subsection*{The exploratory analyses}

The analyses in this subsection were specified on 9 and 10 October 2026, after every registered result was known. They are exploratory, their tests are not corrected as a set, and none changes a registered verdict.

\paragraph*{The pooled model of second works.}
A pooled model fitted the second-work ratings of Studies 2, 4 and 5 together, one model per measure: rating $\sim$ label $*$ modification level $*$ study${}+ (1 \mid \textrm{study}{:}\textrm{participant}) + (1 \mid \textrm{original}) + (1 \mid \textrm{original}{:}\textrm{level})$. It averaged the label's penalty at each level over the three studies with equal weight, whatever their sizes, and averaged Studies 4 and 5 over their time factors, which it leaves out (Supplementary Note~6). It compared four accounts of how the penalty might change with the level: severance (no penalty), taint (the same penalty at every level), dilution (a penalty that changes in a straight line across the levels) and a free shape. Each account was fitted by maximum likelihood, so that their fits can be compared. The comparison used AIC and BIC, two standard scores for choosing between models that reward fit and charge for each added quantity, BIC more heavily. The pooled model treats each study's participants as different people, so anyone who took part in Study~5 and in Study~2 or~4 counts twice. A breakdown of Study~5's registered interaction gave the penalty for each combination of editing time and modification level (Supplementary Note~5).

\paragraph*{The explanations of second works.}
Two analyses described the second-work explanations (Supplementary Note~7): how often they mentioned prompting or machines, by the word lists, and how far each participant's explanations moved from None to Dramatic. The second was measured by word overlap and by a sentence-embedding model (all-mpnet-base-v2), which turns each explanation into a list of numbers whose distances reflect differences in meaning.\cite{reimers2019sentencebert,sentencetransformers2021allmpnet} Each person's distance was taken relative to other participants writing about the same picture, and the labels were compared against 200 shuffles of the label within each study. A structural topic model, which finds groups of words that occur together and relates how often they occur to the conditions, described the pooled explanations (Section~\ref{sec:suppB-crossstudy}).

\paragraph*{The image check.}
An image check measured each edited image's distance from its original with two image-recognition networks, DINOv2 and CLIP,\cite{oquab2024dinov2,radford2021learning} and with a color measure, the difference between the two images' color histograms. It asked whether the versions of each artwork are arranged as the design assumes, by rules fixed in its code before any distance was computed (Supplementary Note~4). Neither network was built to judge art, so neither distance shows how different two images looked to participants.

\subsection*{Software, random seeds and run records}

The models ran in R 4.5.2 (lme4 2.0-1, lmerTest 3.2-1, emmeans 2.0.3, ordinal 2025.12-29, brms 2.23.0 with rstan 2.32.7, glmnet 5.0, stm 1.3.8). The text measures and the image check ran in Python 3.12.3 (pandas 2.2.2, NumPy 1.26.4, scikit-learn 1.5.2, SciPy 1.13.0, sentence-transformers 3.3.1, PyTorch 2.5.1 and, for the image check, transformers 4.46.2). Each run records the versions it used; only the image check also fixes them in a file of requirements. Refitting the four Bayesian models needs a C++ compiler.

Random steps used fixed seeds. A seed is the starting number of a random-number generator, so a rerun with the same seed repeats the same random draws. The R code used 20260421, as did the text code of Studies 2 and 3. Study~1's text code used the text module's defaults, and the text code of Studies 4 and 5 used 20260903, as their registrations name. The image networks were run at fixed versions (DINOv2 revision f9e44c8, CLIP revision 3d74acf). The sentence-embedding model was loaded only from the copy stored on the analysis computer, revision e8c3b32, the only copy there; the run records do not name the revision.

Each analysis writes a run record that lists a checksum, a fingerprint of a file's exact bytes, for every input and output, and later analyses check the files they read against those records. Every result, and every count of participants, records or ratings, printed in the paper is read from the tables that the analysis writes. The script that reads them records the table, row and column behind each number, with the checksum of the table it came from.
